% reguläre Sprache
Über dem Alphabet $\Sigma=\{\texttt{0},\texttt{1}\}$ ist die Sprache
\[
L
=
\{
w\in\Sigma^*
\mid
|w|_{\texttt{0}}
-
|w|_{\texttt{1}}
\not\equiv 0
\mod 3
\}
\]
gegeben.
Sie wird vom endlichen Automaten
\begin{center}
\begin{tikzpicture}[>=latex,thick]
\coordinate (q0) at (0,0);
\coordinate (q1) at (3,{sqrt(3)});
\coordinate (q2) at (3,{-sqrt(3)});
\node at (q0) {$q_0$};
\node at (q1) {$q_1$};
\node at (q2) {$q_2$};
\draw (q0) circle[radius=0.35];
\draw (q1) circle[radius=0.35];
\draw (q1) circle[radius=0.30];
\draw (q2) circle[radius=0.35];
\draw (q2) circle[radius=0.30];
\draw[->,shorten >= 0.35cm] ($(q0)+(-1.5,0)$) -- (q0);
\draw[->,shorten <= 0.35cm,shorten >= 0.35cm] (q0) -- (q1);
\draw[->,shorten <= 0.35cm,shorten >= 0.35cm] (q1) -- (q2);
\draw[->,shorten <= 0.35cm,shorten >= 0.35cm] (q2) -- (q0);
\draw[->,shorten <= 0.35cm,shorten >= 0.35cm] (q0) to[out=-50,in=180] (q2);
\draw[->,shorten <= 0.35cm,shorten >= 0.35cm] (q2) to[out=60,in=-60] (q1);
\draw[->,shorten <= 0.35cm,shorten >= 0.35cm] (q1) to[out=180,in=60] (q0);
\node at ($0.5*(q0)+0.5*(q1)$) [below right] {\texttt{0}};
\node at ($0.5*(q1)+0.5*(q2)$) [left] {\texttt{0}};
\node at ($0.5*(q2)+0.5*(q0)$) [above right] {\texttt{0}};

\node at ($0.65*(q0)+0.65*(q1)-0.3*(q2)$) {\texttt{1}};
\node at ($0.65*(q1)+0.65*(q2)-0.3*(q0)$) {\texttt{1}};
\node at ($0.65*(q2)+0.65*(q0)-0.3*(q1)$) {\texttt{1}};
\end{tikzpicture}
\end{center}
akzeptiert.
Finden Sie einen regulären Ausruck, der $L$ akzeptiert.

\begin{loesung}
\definecolor{darkred}{rgb}{0.8,0,0}
Dazu wird der endliche Automat erst zum folgenden VNEA standardisiert:
\begin{center}
\begin{tikzpicture}[>=latex,thick]
\coordinate (q0) at (0,0);
\coordinate (q1) at (3,{sqrt(3)});
\coordinate (q2) at (3,{-sqrt(3)});
\coordinate (e) at (6,0);
\coordinate (s) at (-3,0);
\node at (s) {$s$};
\node at (q0) {$q_0$};
\node at (q1) {$q_1$};
\node at (q2) {$q_2$};
\node at (e) {$e$};
\draw (s) circle[radius=0.35];
\draw (q0) circle[radius=0.35];
\draw (q1) circle[radius=0.35];
\draw (q2) circle[radius=0.35];
\draw (e) circle[radius=0.30];
\draw (e) circle[radius=0.35];
\draw[->,shorten >= 0.35cm] ($(s)+(-1.5,0)$) -- (s);
\draw[->,shorten <= 0.35cm,shorten >= 0.35cm] (s) -- (q0);
\draw[->,shorten <= 0.35cm,shorten >= 0.35cm] (q0) -- (q1);
\draw[->,shorten <= 0.35cm,shorten >= 0.35cm] (q1) -- (q2);
\draw[->,shorten <= 0.35cm,shorten >= 0.35cm] (q2) -- (q0);
\draw[->,shorten <= 0.35cm,shorten >= 0.35cm] (q1) -- (e);
\draw[->,shorten <= 0.35cm,shorten >= 0.35cm] (q2) -- (e);
\draw[->,shorten <= 0.35cm,shorten >= 0.35cm] (q0) to[out=-50,in=180] (q2);
\draw[->,shorten <= 0.35cm,shorten >= 0.35cm] (q2) to[out=60,in=-60] (q1);
\draw[->,shorten <= 0.35cm,shorten >= 0.35cm] (q1) to[out=180,in=60] (q0);
\node at ($0.5*(q0)+0.5*(q1)$) [below right] {\texttt{0}};
\node at ($0.5*(q1)+0.5*(q2)$) [left] {\texttt{0}};
\node at ($0.5*(q2)+0.5*(q0)$) [above right] {\texttt{0}};
\node at ($0.5*(s)+0.5*(q0)$) [above] {$\varepsilon$};
\node at ($0.65*(q0)+0.65*(q1)-0.3*(q2)$) {\texttt{1}};
\node at ($0.65*(q1)+0.65*(q2)-0.3*(q0)$) {\texttt{1}};
\node at ($0.65*(q2)+0.65*(q0)-0.3*(q1)$) {\texttt{1}};
\node at ($0.5*(q1)+0.5*(e)$) [above right] {$\varepsilon$};
\node at ($0.5*(q2)+0.5*(e)$) [below right] {$\varepsilon$};
\end{tikzpicture}
\end{center}
Jetzt werden die Zwischenzustände in der Reihenfolge $q_2$, $q_1$, $q_0$
entfernt.
Zuerst $q_2$:
\begin{center}
\begin{tikzpicture}[>=latex,thick]
\coordinate (q0) at (0,0);
\coordinate (q1) at (3,{sqrt(3)});
\coordinate (q2) at (3,{-sqrt(3)});
\coordinate (e) at (6,0);
\coordinate (s) at (-3,0);
\node at (s) {$s$};
\node at (q0) {$q_0$};
\node at (q1) {$q_1$};
\node at (e) {$e$};
\draw (s) circle[radius=0.35];
\draw (q0) circle[radius=0.35];
\draw (q1) circle[radius=0.35];
\draw (e) circle[radius=0.30];
\draw (e) circle[radius=0.35];
\draw[->,shorten >= 0.35cm] ($(s)+(-1.5,0)$) -- (s);
\draw[->,shorten <= 0.35cm,shorten >= 0.35cm] (s) -- (q0);
\draw[->,shorten <= 0.35cm,shorten >= 0.35cm] (q0) -- (q1);
\draw[->,shorten <= 0.35cm,shorten >= 0.35cm] (q1) -- (e);
\draw[->,shorten <= 0.35cm,shorten >= 0.35cm] (q1) to[out=180,in=60] (q0);
\draw[->,shorten <= 0.35cm,shorten >= 0.35cm,color=darkred] (q0) -- (e);
\draw[->,shorten <= 0.35cm,shorten >= 0.35cm,color=darkred]
	(q0) to[out=-60,in=-120,distance=1cm] (q0);
\draw[->,shorten <= 0.35cm,shorten >= 0.35cm,color=darkred]
	(q1) to[out=60,in=120,distance=1cm] (q1);
\node[color=darkred] at ($(q1)+(0,1)$) {\texttt{01}};
\node[color=darkred] at ($(q0)+(0,-1)$) {\texttt{10}};
\node at ($0.5*(q0)+0.5*(q1)$) [below right] {\texttt{0{\color{darkred}|11}}};
\node at ($0.5*(s)+0.5*(q0)$) [above] {$\varepsilon$};
\node at ($0.65*(q0)+0.65*(q1)-0.3*(q2)$) {\texttt{1{\color{darkred}|00}}};
\node at ($0.5*(q1)+0.5*(e)$) [above right] {$\varepsilon$\texttt{\color{darkred}|0}};
\node[color=darkred] at ($0.5*(q0)+0.5*(e)$) [below] {$\texttt{1}$};
\end{tikzpicture}
\end{center}
Dann $q_1$:
\begin{center}
\begin{tikzpicture}[>=latex,thick]
\coordinate (q0) at (0,0);
\coordinate (e) at (6,0);
\coordinate (s) at (-3,0);
\node at (s) {$s$};
\node at (q0) {$q_0$};
\node at (e) {$e$};
\draw (s) circle[radius=0.35];
\draw (q0) circle[radius=0.35];
\draw (e) circle[radius=0.30];
\draw (e) circle[radius=0.35];
\draw[->,shorten >= 0.35cm] ($(s)+(-1.5,0)$) -- (s);
\draw[->,shorten <= 0.35cm,shorten >= 0.35cm] (s) -- (q0);
\draw[->,shorten <= 0.35cm,shorten >= 0.35cm] (q0) -- (e);
\draw[->,shorten <= 0.35cm,shorten >= 0.35cm]
	(q0) to[out=-60,in=-120,distance=1cm] (q0);
\node at ($(q0)+(0,-1)$)
	{\texttt{10{\color{darkred}|((0|11)(01)*(1|00))}}};
\node at ($0.5*(s)+0.5*(q0)$) [above] {$\varepsilon$};
\node at ($0.5*(q0)+0.5*(e)$) [below]
	{$\texttt{1{\color{darkred}|(0|11)(01)*($\varepsilon$|0)}}$};
\end{tikzpicture}
\end{center}
Im letzten Schritt wird $q_0$ entfernt, jetzt kann man den regulären
Ausdruck direkt ablesen, er ist
\[
r
=
\texttt{(10|((0|11)(01)*(1|00)))*(1|(0|11)(01)*(|0))}.
\qedhere
\]
\end{loesung}

\begin{diskussion}
\definecolor{darkred}{rgb}{0.8,0,0}
Natürlich kann man auch eine andere Reihenfolge für die Entfernung der
Zwischenzustände wählen.
Beginnt man mit $q_0$, entsteht nach dem ersten Schritt der Automat
\begin{center}
\begin{tikzpicture}[>=latex,thick]
\coordinate (q0) at (0,0);
\coordinate (q1) at (3,{sqrt(3)});
\coordinate (q2) at (3,{-sqrt(3)});
\coordinate (e) at (6,0);
\coordinate (s) at (-3,0);
\node at (s) {$s$};
\node at (q1) {$q_1$};
\node at (q2) {$q_2$};
\node at (e) {$e$};
\draw (s) circle[radius=0.35];
\draw (q1) circle[radius=0.35];
\draw (q2) circle[radius=0.35];
\draw (e) circle[radius=0.30];
\draw (e) circle[radius=0.35];
\draw[->,shorten >= 0.35cm] ($(s)+(-1.5,0)$) -- (s);
\draw[->,shorten <= 0.35cm,shorten >= 0.35cm,color=darkred] (s) -- (q1);
\draw[->,shorten <= 0.35cm,shorten >= 0.35cm,color=darkred] (s) -- (q2);
\draw[->,shorten <= 0.35cm,shorten >= 0.35cm] (q1) -- (q2);
\draw[->,shorten <= 0.35cm,shorten >= 0.35cm] (q1) -- (e);
\draw[->,shorten <= 0.35cm,shorten >= 0.35cm] (q2) -- (e);
\draw[->,shorten <= 0.35cm,shorten >= 0.35cm] (q2) to[out=60,in=-60] (q1);
\node[color=darkred] at ($0.5*(s)+0.5*(q1)$) [above left] {\texttt{0}};
\node[color=darkred] at ($0.5*(s)+0.5*(q2)$) [below left] {\texttt{1}};
\node at ($0.5*(q1)+0.5*(q2)$) [left]
	{\texttt{0}{\color{darkred}\texttt{|11}}};
\node at ($0.60*(q1)+0.60*(q2)-0.3*(q0)$) [right]
	{\texttt{1}{\color{darkred}\texttt{|00}}};
\node at ($0.5*(q1)+0.5*(e)$) [above right] {$\varepsilon$};
\node at ($0.5*(q2)+0.5*(e)$) [below right] {$\varepsilon$};
\draw[->,shorten <= 0.35cm,shorten >= 0.35cm,color=darkred]
	(q1) to[out=60,in=120,distance=1cm] (q1);
\node[color=darkred] at ($(q1)+(0,1)$) {\texttt{10}};
\draw[->,shorten <= 0.35cm,shorten >= 0.35cm,color=darkred]
	(q2) to[out=-60,in=-120,distance=1cm] (q2);
\node[color=darkred] at ($(q2)+(0,-1)$) {\texttt{01}};
\end{tikzpicture}
\end{center}
Entfernung von $q_1$:
\begin{center}
\begin{tikzpicture}[>=latex,thick]
\coordinate (q0) at (0,0);
\coordinate (q1) at (3,{sqrt(3)});
\coordinate (q2) at (3,{-sqrt(3)});
\coordinate (e) at (6,0);
\coordinate (s) at (-3,0);
\node at (s) {$s$};
\node at (q2) {$q_2$};
\node at (e) {$e$};
\draw (s) circle[radius=0.35];
\draw (q2) circle[radius=0.35];
\draw (e) circle[radius=0.30];
\draw (e) circle[radius=0.35];
\draw[->,shorten >= 0.35cm] ($(s)+(-1.5,0)$) -- (s);
\draw[->,shorten <= 0.35cm,shorten >= 0.35cm,color=darkred] (s) -- (e);
\draw[->,shorten <= 0.35cm,shorten >= 0.35cm] (s) -- (q2);
\draw[->,shorten <= 0.35cm,shorten >= 0.35cm] (q2) -- (e);
\node[color=darkred] at ($0.5*(s)+0.5*(e)$) [above]
	{\texttt{0(10)*}};
\node at ($0.3*(s)+0.7*(q2)$) [below left]
	{\texttt{1}\color{darkred}\texttt{|0(10)*(0|11)}};
\node at ($0.5*(q2)+0.5*(e)$) [below right]
	{$\varepsilon${\color{darkred}\texttt{|(1|00)(10)*)}}};
\draw[->,shorten <= 0.35cm,shorten >= 0.35cm,color=darkred]
	(q2) to[out=-60,in=-120,distance=1cm] (q2);
\node at ($(q2)+(0,-1)$)
	{\texttt{01}{\color{darkred}\texttt{|(1|00)(10)*(0|11)}}};
\end{tikzpicture}
\end{center}
Im letzten Schritt, der Entfernung von $q_2$, kann man den resultierenden
regulären Ausdruck direkt hinschreiben:
\[
r=
\texttt{0(10)*|((1|0(10)*(0|11))(01|(1|00)(10)*(0|11))*(|(1|00)(10)*))}
\]
\end{diskussion}

\begin{bewertung}
Neuer Startzustand ({\bf S}) 1 Punkt,
neuer Akzeptierzustand ({\bf A}) 1 Punkt,
Idee Zwischenzustände entfernen ({\bf E}) 1 Punkt,
Regel für Entfernung eines Zwischenzustands ({\bf Z}) 1 Punkt,
Entstehung neuer Schleifen ({\bf L}) 1 Punkt,
regulärer Ausdruck ({\bf R}) 1 Punkt.
\end{bewertung}

